

\section{Цель работы}

Целью данной лабораторной работы является освоение метода анализа иерархий (МАИ) для принятия решений в условиях неопределенности.

В рамках работы требуется освоить методику парных сравнений, которая позволяет переводить качественные суждения экспертов в количественные оценки. Необходимо понять, как вычисляются локальные приоритеты элементов на каждом уровне иерархии на основе собственных векторов матриц парных сравнений. Важно усвоить процедуру проверки согласованности матриц, которая гарантирует логическую непротиворечивость экспертных оценок. Целью также является приобретение навыков синтеза глобальных приоритетов альтернатив путем взвешенного суммирования локальных приоритетов по всем уровням иерархии. Кроме того, работа направлена на развитие умения интерпретировать полученные результаты и выбирать оптимальную альтернативу на основе рассчитанных приоритетов.

В конечном счете, выполнение работы должно сформировать практические навыки применения МАИ для решения реальных задач многокритериального выбора. Также целью является закрепление теоретических знаний, полученных на лекциях, и их применение к конкретному варианту задания.

\clearpage

\section{Теоретическая часть}

Метод анализа иерархий (МАИ), разработанный Томасом Саати, является систематической процедурой для иерархического представления элементов, определяющих суть проблемы. Метод состоит в декомпозиции проблемы на более простые составляющие части и дальнейшей обработке последовательности суждений лица, принимающего решения, по парным сравнениям. В результате может быть выражена относительная степень взаимодействия элементов в иерархии. Эти суждения затем выражаются численно.

В общем случае иерархия строится с вершины (цели), через промежуточные уровни (критерии) к самому нижнему уровню (альтернативы). После построения иерархии устанавливаются приоритеты критериев и оцениваются каждое из альтернатив по каждому критерию. Эти приоритеты выражаются числами. Основным инструментом МАИ является матрица парных сравнений, элементы которой представляют собой оценки превосходства одного элемента над другим по отношению к некоторому критерию.


Матрица парных сравнений $A = (a_{ij})$ является положительной обратно-симметричной матрицей, для которой выполнены условия:
\begin{itemize}
    \item $a_{ij} > 0$;
    \item $a_{ij} = 1/a_{ji}$ (обратная симметричность);
    \item $a_{ii} = 1$ (диагональные элементы равны единице).
\end{itemize}

Элементы $a_{ij}$ отражают степень превосходства элемента $i$ над элементом $j$. Для заполнения матрицы используется шкала относительной важности Саати, которая включает следующие значения: 1 -- равная важность; 3 -- умеренное превосходство; 5 -- существенное превосходство; 7 -- значительное превосходство; 9 -- очень сильное превосходство. Значения 2, 4, 6, 8 обычно используются для промежуточных решений.


Для каждой матрицы парных сравнений необходимо вычислить вектор локальных приоритетов $w = (w_1, w_2, \ldots, w_n)^T$. В методе анализа иерархий используется подход, основанный на вычислении главного собственного вектора матрицы $A$, соответствующего максимальному собственному значению $\lambda_{\max}$:
\[
A w = \lambda_{\max} w.
\]

На практике для приближенного вычисления вектора приоритетов часто используется метод геометрического среднего. Сначала вычисляется геометрическое среднее каждой строки матрицы:
\[
\tilde{w}_i = \sqrt[n]{\prod_{j=1}^{n} a_{ij}}, \quad i = 1, \ldots, n.
\]

Затем полученные значения нормализуются:
\[
w_i = \frac{\tilde{w}_i}{\sum_{k=1}^{n} \tilde{w}_k}, \quad i = 1, \ldots, n.
\]

Полученный вектор $w = (w_1, w_2, \ldots, w_n)^T$ является вектором локальных приоритетов, где $w_i$ -- вес (важность) $i$-го элемента, причем $\sum_{i=1}^{n} w_i = 1$.



После вычисления локальных приоритетов для всех матриц парных сравнений (как для критериев, так и для альтернатив по каждому критерию) выполняется синтез глобальных приоритетов альтернатив. Глобальный приоритет $i$-й альтернативы вычисляется по формуле:
\[
W_i = \sum_{k=1}^{m} e_k \cdot w_{i,k}^{(alt)},
\]
где $m$ -- количество критериев; $e_k$ -- вес $k$-го критерия; $w_{i,k}^{(alt)}$ -- локальный приоритет $i$-й альтернативы по $k$-му критерию.

Оптимальной считается альтернатива, имеющая максимальный глобальный приоритет.


\clearpage
\section{Практическая часть}

В соответствии с вариантом задания необходимо найти приоритеты альтернатив и определить оптимальную альтернативу методом анализа иерархий. Исходными данными являются матрица парных сравнений критериев $E$ (размерность $7 \times 7$) и матрицы парных сравнений альтернатив по каждому из семи критериев $W_1, W_2, \ldots, W_7$ (размерность $6 \times 6$).


На первом этапе необходимо вычислить вектор приоритетов критериев на основе матрицы $E$. Для этого используем метод геометрического среднего.

Матрица $E$ имеет вид:

\[
E = \begin{pmatrix}
1 & 1/3 & 4 & 1/3 & 1/5 & 4 & 8 \\
3 & 1 & 1/3 & 1/7 & 1/2 & 1/6 & 9 \\
1/4 & 3 & 1 & 1/8 & 2 & 1/5 & 9 \\
3 & 7 & 8 & 1 & 9 & 1/6 & 1/7 \\
5 & 2 & 1/2 & 1/9 & 1 & 3 & 1/2 \\
1/4 & 6 & 5 & 6 & 1/3 & 1 & 1/2 \\
1/8 & 1/9 & 1/9 & 7 & 2 & 2 & 1
\end{pmatrix}
\]

Вычислим геометрическое среднее для каждой строки матрицы $E$:

\begin{align*}
\tilde{e}_1 &= \sqrt[7]{1 \cdot \frac{1}{3} \cdot 4 \cdot \frac{1}{3} \cdot \frac{1}{5} \cdot 4 \cdot 8} = 1.161065657, \\
\tilde{e}_2 &= \sqrt[7]{3 \cdot 1 \cdot \frac{1}{3} \cdot \frac{1}{7} \cdot \frac{1}{2} \cdot \frac{1}{6} \cdot 9} = 0.726814056, \\
\tilde{e}_3 &= \sqrt[7]{\frac{1}{4} \cdot 3 \cdot 1 \cdot \frac{1}{8} \cdot 2 \cdot \frac{1}{5} \cdot 9} = 0.8562696276, \\
\tilde{e}_4 &= \sqrt[7]{3 \cdot 7 \cdot 8 \cdot 1 \cdot 9 \cdot \frac{1}{6} \cdot \frac{1}{7}} = 1.668510441, \\
\tilde{e}_5 &= \sqrt[7]{5 \cdot 2 \cdot \frac{1}{2} \cdot \frac{1}{9} \cdot 1 \cdot 3 \cdot \frac{1}{2}} = 0.9742903329, \\
\tilde{e}_6 &= \sqrt[7]{\frac{1}{4} \cdot 6 \cdot 5 \cdot 6 \cdot \frac{1}{3} \cdot 1 \cdot \frac{1}{2}} = 1.333548306, \\
\tilde{e}_7 &= \sqrt[7]{\frac{1}{8} \cdot \frac{1}{9} \cdot \frac{1}{9} \cdot 7 \cdot 2 \cdot 2 \cdot 1} = 0.6383859046.
\end{align*}

Сумма полученных значений равна $7.358903326$. После нормализации получаем вектор приоритетов критериев:
\begin{gather*}
e = (0.1577774029,\ 0.09876688149,\ 0.1163586204,\ 0.226734158,\\ 0.1323964734,\ 0.1812160984,\ 0.08675036547)^T.
\end{gather*}


Теперь вычислим локальные приоритеты альтернатив для каждой из семи матриц $W_1, \ldots, W_7$. Для каждой матрицы вычислим геометрические средние строк, а затем выполним их нормализацию.


Матрица $W_1$:
\[
W_1 = \begin{pmatrix}
1 & 5 & 1/8 & 1/3 & 1/3 & 1/5 \\
1/5 & 1 & 8 & 1/4 & 1/4 & 4 \\
8 & 1/8 & 1 & 1/6 & 2 & 8 \\
3 & 4 & 6 & 1 & 9 & 8 \\
3 & 4 & 1/2 & 1/9 & 1 & 8 \\
5 & 1/4 & 1/8 & 1/8 & 1/8 & 1
\end{pmatrix}
\]

Геометрические средние строк до нормализации:
\begin{gather*}
\tilde{w}^1 = (0.4902804589,\ 0.8583742189,\ 1.177591843,\ 4.160167646,\\ 1.321802152,\ 0.3669498169)^T.
\end{gather*}

Сумма значений равна $8.375166136$. После нормализации получаем вектор локальных приоритетов:
\begin{gather*}
w^1 = (0.05853978906,\ 0.1024904109,\ 0.140605192,\ 0.4967265817,\\ 0.1578239919,\ 0.04381403437)^T.
\end{gather*}



Матрица $W_2$:
\[
W_2 = \begin{pmatrix}
1 & 6 & 1/5 & 1/8 & 2 & 6 \\
1/6 & 1 & 3 & 3 & 1/2 & 1/3 \\
5 & 1/3 & 1 & 1/6 & 5 & 7 \\
8 & 1/3 & 6 & 1 & 3 & 4 \\
1/2 & 2 & 1/5 & 1/3 & 1 & 1/7 \\
1/6 & 3 & 1/7 & 1/4 & 7 & 1
\end{pmatrix}
\]

Геометрические средние строк до нормализации:
\begin{gather*}
\tilde{w}^2 = (1.102923569,\ 0.793700526,\ 1.460923889,\ 2.40187391,\\ 0.46039979,\ 0.7071067812)^T.
\end{gather*}

Сумма значений равна $6.926928466$. После нормализации получаем вектор локальных приоритетов:
\begin{gather*}
w^2 = (0.1592226013,\ 0.1145818857,\ 0.2109050059,\ 0.3467444369,\\ 0.06646521503,\ 0.1020808551)^T.
\end{gather*}


Матрица $W_3$:
\[
W_3 = \begin{pmatrix}
1 & 1/5 & 8 & 1/7 & 1 & 4 \\
5 & 1 & 8 & 1/5 & 1/7 & 1/5 \\
1/8 & 1/8 & 1 & 1/5 & 7 & 1/6 \\
7 & 5 & 5 & 1 & 1/2 & 1/2 \\
1 & 7 & 1/7 & 2 & 1 & 1/3 \\
1/4 & 5 & 6 & 2 & 3 & 1
\end{pmatrix}
\]

Геометрические средние строк до нормализации:
\begin{gather*}
\tilde{w}^3 = (0.9851756195,\ 0.7819344074,\ 0.3923130925,\ 1.877138611,\\ 0.9346552652,\ 1.885972774)^T.
\end{gather*}

Сумма значений равна $6.85718977$. После нормализации получаем вектор локальных приоритетов:
\begin{gather*}
w^3 = (0.1436704616,\ 0.1140313209,\ 0.05721193458,\ 0.273747508,\\ 0.1363029603,\ 0.2750358146)^T.
\end{gather*}


Матрица $W_4$:
\[
W_4 = \begin{pmatrix}
1 & 5 & 7 & 3 & 1/3 & 1/6 \\
1/5 & 1 & 8 & 1/4 & 1/7 & 1/9 \\
1/7 & 1/8 & 1 & 2 & 9 & 6 \\
1/3 & 4 & 1/2 & 1 & 2 & 1/9 \\
3 & 7 & 1/9 & 1/2 & 1 & 1/9 \\
6 & 9 & 1/6 & 9 & 9 & 1
\end{pmatrix}
\]

Геометрические средние строк до нормализации:
\begin{gather*}
\tilde{w}^4 = (1.341691903,\ 0.4303150878,\ 1.115679076,\ 0.7274157573,\\ 0.7114057647,\ 3)^T.
\end{gather*}

Сумма значений равна $7.326507589$. После нормализации получаем вектор локальных приоритетов:
\begin{gather*}
w^4 = (0.1831284397,\ 0.05873399878,\ 0.1522797953,\ 0.09928547108,\\ 0.0971002563,\ 0.4094720388)^T.
\end{gather*}


Матрица $W_5$:
\[
W_5 = \begin{pmatrix}
1 & 9 & 1/7 & 1/7 & 7 & 1/3 \\
1/9 & 1 & 4 & 3 & 8 & 3 \\
7 & 1/4 & 1 & 9 & 1/5 & 4 \\
7 & 1/3 & 1/9 & 1 & 6 & 1/9 \\
1/7 & 1/8 & 5 & 1/6 & 1 & 6 \\
3 & 1/3 & 1/4 & 9 & 1/6 & 1
\end{pmatrix}
\]

Геометрические средние строк до нормализации:
\begin{gather*}
\tilde{w}^5 = (0.868301469,\ 1.781797436,\ 1.525439862,\ 0.7463465039,\\ 0.6685445143,\ 0.8491906648)^T.
\end{gather*}

Сумма значений равна $6.439620451$. После нормализации получаем вектор локальных приоритетов:
\begin{gather*}
w^5 = (0.1348373675,\ 0.2766929278,\ 0.2368835048,\ 0.1158991449,\\ 0.1038173786,\ 0.1318696764)^T.
\end{gather*}


Матрица $W_6$:
\[
W_6 = \begin{pmatrix}
1 & 1/3 & 3 & 3 & 3 & 1/4 \\
3 & 1 & 1/4 & 1/7 & 7 & 6 \\
1/3 & 4 & 1 & 1/3 & 6 & 1/2 \\
1/3 & 7 & 3 & 1 & 1/5 & 2 \\
1/3 & 1/7 & 1/6 & 5 & 1 & 1/6 \\
4 & 1/6 & 2 & 1/2 & 6 & 1
\end{pmatrix}
\]

Геометрические средние строк до нормализации:
\begin{gather*}
\tilde{w}^6 = (1.144714243,\ 1.284898293,\ 1.049115063,\ 1.187206699,\\ 0.43325279,\ 1.25992105)^T.
\end{gather*}

Сумма значений равна $6.359108138$. После нормализации получаем вектор локальных приоритетов:
\begin{gather*}
w^6 = (0.1800117591,\ 0.2020563679,\ 0.1649783335,\ 0.1866939,\\ 0.06813106187,\ 0.1981285775)^T.
\end{gather*}


Матрица $W_7$:
\[
W_7 = \begin{pmatrix}
1 & 1/9 & 1/3 & 9 & 7 & 7 \\
9 & 1 & 1/8 & 1/5 & 7 & 8 \\
3 & 8 & 1 & 9 & 6 & 1/3 \\
1/9 & 5 & 1/9 & 1 & 6 & 9 \\
1/7 & 1/7 & 1/6 & 1/6 & 1 & 9 \\
1/7 & 1/8 & 3 & 1/9 & 1/9 & 1
\end{pmatrix}
\]

Геометрические средние строк до нормализации:
\begin{gather*}
\tilde{w}^7 = (1.592865616,\ 1.525439862,\ 2.749459274,\ 1.222211758,\\ 0.4149132667,\ 0.2951716897)^T.
\end{gather*}

Сумма значений равна $7.800062067$. После нормализации получаем вектор локальных приоритетов:
\begin{gather*}
w^7 = (0.2042119313,\ 0.1955676719,\ 0.352492001,\ 0.1566925804,\\ 0.05319358937,\ 0.03784222612)^T.
\end{gather*}


Теперь вычислим глобальные приоритеты альтернатив по формуле:
\[
W_i = \sum_{k=1}^{7} e_k \cdot w_i^k, \quad i = 1, \ldots, 6.
\]

Подставим полученные значения:

\[
W_1 = 0.1577774029 \cdot 0.05853978906 + 0.09876688149 \cdot 0.1592226013 +
\]
\[
+ 0.1163586204 \cdot 0.1436704616 + 0.226734158 \cdot 0.1831284397 +
\]
\[
+ 0.1323964734 \cdot 0.1348373675 + 0.1812160984 \cdot 0.1800117591 +
\]
\[
+ 0.08675036547 \cdot 0.2042119313 =
\]
\[
= 0.0092 + 0.0157 + 0.0167 + 0.0415 + 0.0178 + 0.0326 + 0.0177 = 0.1513894252.
\]

\[
W_2 = 0.1577774029 \cdot 0.1024904109 + 0.09876688149 \cdot 0.1145818857 +
\]
\[
+ 0.1163586204 \cdot 0.1140313209 + 0.226734158 \cdot 0.05873399878 +
\]
\[
+ 0.1323964734 \cdot 0.2766929278 + 0.1812160984 \cdot 0.2020563679 +
\]
\[
+ 0.08675036547 \cdot 0.1955676719 =
\]
\[
= 0.0162 + 0.0113 + 0.0133 + 0.0133 + 0.0366 + 0.0366 + 0.0170 = 0.1442876988.
\]

\[
W_3 = 0.1577774029 \cdot 0.140605192 + 0.09876688149 \cdot 0.2109050059 +
\]
\[
+ 0.1163586204 \cdot 0.05721193458 + 0.226734158 \cdot 0.1522797953 +
\]
\[
+ 0.1323964734 \cdot 0.2368835048 + 0.1812160984 \cdot 0.1649783335 +
\]
\[
+ 0.08675036547 \cdot 0.352492001 =
\]
\[
= 0.0222 + 0.0208 + 0.0067 + 0.0345 + 0.0314 + 0.0299 + 0.0306 = 0.1760369652.
\]

\[
W_4 = 0.1577774029 \cdot 0.4967265817 + 0.09876688149 \cdot 0.3467444369 +
\]
\[
+ 0.1163586204 \cdot 0.273747508 + 0.226734158 \cdot 0.09928547108 +
\]
\[
+ 0.1323964734 \cdot 0.1158991449 + 0.1812160984 \cdot 0.1866939 +
\]
\[
+ 0.08675036547 \cdot 0.1566925804 =
\]
\[
= 0.0784 + 0.0343 + 0.0319 + 0.0225 + 0.0153 + 0.0338 + 0.0136 = 0.2297531036.
\]

\[
W_5 = 0.1577774029 \cdot 0.1578239919 + 0.09876688149 \cdot 0.06646521503 +
\]
\[
+ 0.1163586204 \cdot 0.1363029603 + 0.226734158 \cdot 0.0971002563 +
\]
\[
+ 0.1323964734 \cdot 0.1038173786 + 0.1812160984 \cdot 0.06813106187 +
\]
\[
+ 0.08675036547 \cdot 0.05319358937 =
\]
\[
= 0.0249 + 0.0066 + 0.0159 + 0.0220 + 0.0137 + 0.0123 + 0.0046 = 0.1000476542.
\]

\[
W_6 = 0.1577774029 \cdot 0.04381403437 + 0.09876688149 \cdot 0.1020808551 +
\]
\[
+ 0.1163586204 \cdot 0.2750358146 + 0.226734158 \cdot 0.4094720388 +
\]
\[
+ 0.1323964734 \cdot 0.1318696764 + 0.1812160984 \cdot 0.1981285775 +
\]
\[
+ 0.08675036547 \cdot 0.03784222612 =
\]
\[
= 0.0069 + 0.0101 + 0.0320 + 0.0928 + 0.0175 + 0.0359 + 0.0033 = 0.198485153.
\]

Результаты вычислений глобальных приоритетов:

\[
\begin{gathered}
x_1 = 0.1513894252 \\
x_2 = 0.1442876988 \\
x_3 = 0.1760369652 \\
x_4 = \mathbf{0.2297531036} \\
x_5 = 0.1000476542 \\
x_6 = 0.198485153
\end{gathered}
\]
Максимальный глобальный приоритет имеет альтернатива $x_4$ со значением $W_4 = 0.2297531036$. Следовательно, оптимальной альтернативой является $x_4$.


\clearpage
\section{Заключение}

В ходе выполнения данной лабораторной работы был освоен метод анализа иерархий для принятия решений в условиях неопределенности. Была построена иерархическая структура задачи, включающая цель, семь критериев и шесть альтернатив. Для каждого уровня иерархии были составлены матрицы парных сравнений и вычислены локальные приоритеты элементов. С использованием метода геометрического среднего были получены векторы приоритетов для матрицы критериев $E$ и для семи матриц альтернатив $W_1, \ldots, W_7$. Проведен синтез глобальных приоритетов альтернатив путем взвешенного суммирования локальных приоритетов с учетом весов критериев. В результате вычислений было установлено, что оптимальной альтернативой является $x_4$, имеющая максимальный глобальный приоритет $W_4 = 0.229753$.

Полученные результаты демонстрируют эффективность применения метода анализа иерархий для решения задач многокритериального выбора. Выполнение работы позволило закрепить теоретические знания и приобрести практические навыки использования МАИ. Таким образом, цель лабораторной работы была полностью достигнута.

\clearpage
\section{Список литературы}

\begin{enumerate}
    \item Саати Т. Принятие решений при зависимостях и обратных связях: Аналитические сети. --- М.: Книжный дом «ЛИБРОКОМ», 2009. --- 360 с.
    \item Теория принятия решений в 2 т. Том 2: учебник и практикум для вузов / В. Г. Халин [и др.]; ответственный редактор В. Г. Халин. --- Москва: Издательство Юрайт, 2020. --- 431 с.
    \item Баранов, В. В. Процессы принятия управляющих решений, мотивированных интересами: учебное пособие / В. В. Баранов. --- Москва: ФИЗМАТЛИТ, 2005. --- 296 с.
    \item Петросян Л.А., Зенкевич Н.А., Шевкопляс Е.В. Теория игр. --- СПб.: БХВ-Петербург, 2012. --- 432 с.
    \item Лубенец Ю.В. Теория принятия решений: Учеб. пособие. --- Липецк: Изд-во ЛГТУ, 2011. --- 80 с.
\end{enumerate}
